% !PRINTED 197
\weberpage{197}

Die Funktion (1) verschwindet für alle von Null verschiedenen Werte $u^0$
von $u$, die der Bedingung
\[
\pm n\,u^0 \equiv u^0 \pmod{\omega_1,\omega_2}
\]
genügen, also für
\weq{4}{u^0 = \frac{\nu_1\omega_1 + \nu_2\omega_2}{n\pm 1},}
wenn wieder $\nu_1, \nu_2$ beliebige ganze Zahlen sind. Die Nullpunkte sind
von der ersten, die Unendlichkeitspunkte von der zweiten Ordnung.

Wir führen daher jetzt eine Funktion $\psi_n(u)$ ein, die wir
folgendermaßen erklären:
\weq{5}{\psi_n(u)^2 = n^2 \prod^{\nu_1,\nu_2}\left[\wp(u) - \wp\left(\frac{\nu_1\omega_1+\nu_2\omega_2}{n}\right)\right],}
wobei $\nu_1, \nu_2$ je ein vollständiges Restsystem nach dem Modul $n$,
mit alleinigem Ausschluß des Wertepaares $0,0$ durchlaufen, und fügen noch
die Bestimmung hinzu, daß $\psi_1 = 1$ sein soll.

Wenn $n$ ungerade ist, so kommt in dem Produkt (5) jeder Linearfaktor
zweimal vor, da die beiden inkongruenten Werte
\[
u = \pm\frac{\nu_1\omega_1 + \nu_2\omega_2}{n}
\]
den gleichen Wert von $\wp(u)$ ergeben. Ist aber $n$ gerade, so kommen
wieder in (5) alle Linearfaktoren zweimal vor, mit Ausnahme der drei
\[
\wp(u) - \wp\left(\frac{\omega_1}{2}\right),\quad
\wp(u) - \wp\left(\frac{\omega_2}{2}\right),\quad
\wp(u) - \wp\left(\frac{\omega_1+\omega_2}{2}\right),
\]
die nur einfach vorkommen. Beachtet man nun, daß nach § 41
\[
\wp'(u)^2 = 4\left[\wp(u)-\wp\left(\frac{\omega_1}{2}\right)\right]\left[\wp(u)-\wp\left(\frac{\omega_2}{2}\right)\right]\left[\wp(u)-\wp\left(\frac{\omega_1+\omega_2}{2}\right)\right]
\]
ist, so ergibt sich das Resultat:
\weq{6}{
\begin{array}{ll}
n\text{ ungerade:} & \psi_n = P_n,\\
n\text{ gerade:} & \psi_n = \wp'(u)\,P_n,
\end{array}}
wenn $P_n$ eine ganze rationale Funktion von $\wp(u)$ bedeutet, die bei
ungeradem $n$ vom Grade $\frac12(n^2-1)$ und bei geradem $n$ vom Grade
$\frac12(n^2-4)$ ist, und in der das Glied höchster Ordnung bzw. gleich
\[
n\,\wp(u)^{\frac{n^2-1}{2}},\qquad -\frac{n}{2}\,\wp(u)^{\frac{n^2-4}{2}}
\]
%CONT>>
